\section{Họ}

Cho một tập hợp $x$. \defText{Họ phần tử} của $x$ (gọi tắt là \defText{họ}) là một ánh xạ bất kì có tập dích là $x$. Nếu ánh xạ này là $\begin{array}[t]{rcl}  
f: I & \to & x\\       
   a & \mapsto & f(a) 
\end{array}$ thì chúng ta kí hiệu họ phần tử này là $\defMath{(f_i)_{i \in I}}$, và tập nguồn $I$ được gọi là \defText{tập chỉ số} của họ phần tử đó.

Dưới đây là định nghĩa của một số các cách viết quen thuộc. Cho $x$ là một tập hợp. Một họ $(f_i)_{i \in I}$ được xác định trong đó $\forall i \in I, f_i \in \mathfrak{P}(x)$. Khi đó, \begin{itemize}
   \item \defText{hợp của họ}, kí hiệu là $\defMath{\bigcup_{i \in I}f_i}$, được định nghĩa bởi $$
   \defMath{\bigcup_{i \in I}f_i = \{y \in x \mid \exists i \in I, y \in f_i\}};
   $$
   \item \defText{giao của họ}, kí hiệu là $\defMath{\bigcap_{i \in I}f_i}$, được định nghĩa bởi $$
   \defMath{\bigcap_{i \in I}f_i = \{y \in x \mid \forall i \in I, y \in f_i\}}.$$
\end{itemize}
Như thế, theo định nghĩa này, với mọi $y \in x$, chúng ta có khẳng định
\begin{multicols}{2}
   \begin{itemize}
      \item $y \in \bigcup_{i \in I}f_i \iff \exists i \in I, y \in f_i$;
      \item $y \in \bigcap_{i \in I}f_i \iff \forall i \in I, y \in f_i$.
   \end{itemize}
\end{multicols}
Ngoài ra, nếu xét tập chỉ số $I$ là một tập rỗng, thì chúng ta có $i \in I$ luôn sai với mọi $i$, do đó 
\begin{multicols}{2}
   \begin{itemize}
      \item $\bigcup_{i \in \varnothing}f_i = \varnothing$;
      \item $\bigcap_{i \in \varnothing}f_i = x$.
   \end{itemize}
\end{multicols}

\exercise Cho tập hợp $x$, $M$ là một tập con của $x$, $A$ và $B$ đều là ánh xạ từ $I$ đến $\mathfrak{P}(x)$. Chứng minh rằng
\begin{multicols}{2}
   \begin{enumerate}
      \item $x \setminus \left(\bigcup_{i \in I}A_i\right) = \bigcap_{i \in I}(x \setminus A_i)$;
      \item $x \setminus \left(\bigcap_{i \in I}A_i\right) = \bigcup_{i \in I}(x \setminus A_i)$;
      \item $\left(\bigcap_{i \in I}A_i\right) \cap M = \bigcap_{i \in I}\left(A_i \cap M\right)$;
      \item $I \neq \varnothing \implies \left(\bigcup_{i \in I}A_i\right) \cup M = \bigcup_{i \in I}\left(A_i \cup M\right)$;
      \item $\left(\bigcup_{i \in I}A_i\right) \cap M = \bigcup_{i \in I}\left(A_i \cap M\right)$;
      \item $\left(\bigcap_{i \in I}A_i\right) \cup M = \bigcap_{i \in I}\left(A_i \cup M\right)$;
      \item $\forall i \in I, A_i \subseteq B_i \implies \begin{cases}
         \bigcup_{i \in I} A_i \subseteq \bigcup_{i \in I}B_i \\
         \bigcap_{i \in I} A_i \subseteq \bigcap_{i \in I}B_i \\
      \end{cases}$;
      \item $\bigcap_{i \in I}\left(A_i \setminus B_i\right) = \left(\bigcap_{i \in I}A_i\right) \setminus \left(\bigcup_{i \in I}B_i\right)$.
   \end{enumerate} 
\end{multicols}

\solution 

\resetSubexercise Với mỗi $i$ trong $I$, $A_i$ được xác định duy nhất, dẫn đến $x \setminus A_i$ cũng xác định duy nhất. Ngoài ra, có $\forall i \in I, x \setminus A_i \subseteq x$ nên $\left(x \setminus A_i\right)_{i \in I}$ là một họ phần tử của $x$.

Trong trường hợp $I = \varnothing$, có ngay $\bigcup_{i \in I}A_i = \varnothing$ và $\bigcap_{i \in I}A_i = x$. Vì $x \setminus \varnothing = x$ nên chúng ta có điều phải chứng minh trong trường hợp này. Còn nếu $I \neq \varnothing$, với mọi $e$, chúng ta có 
\begin{equation*}
   e \in x \setminus \left(\bigcup_{i \in I}A_i\right) \iff e \in x \land e \notin \bigcup_{i \in I}A_i.
\end{equation*}
$e \in x \implies \forall i \in I, e \in x$ là hiển nhiên. Xét chiều ngược lại, giả sử $\forall i \in I, e \in x$. Do đang xét $I \neq \varnothing$, tồn tại $j \in I$, dẫn đến khi có $j \in I \implies e \in x$ từ cụ thể hóa lượng từ phổ quát, chúng ta suy ra được $e \in x$. Vì vậy, \begin{equation*}
   e \in x \iff \forall i \in I, e \in x.
\end{equation*} 
Ngoài ra, do $e \in x$, chúng ta cũng có 
\begin{align*}
   e \notin \bigcup_{i \in I}A_i &\iff \overline{\exists i \in I, e \in A_i}; \\
   e \notin \bigcup_{i \in I}A_i &\iff \forall i \in I, e \notin A_i.
\end{align*}

Do đó, \begin{align*}
   e \in x \setminus \left(\bigcup_{i \in I}A_i\right) &\iff \forall i \in I, e \in x \land \forall i \in I, e \notin A_i; \\
   e \in x \setminus \left(\bigcup_{i \in I}A_i\right) &\iff \forall i \in I, \left(e \in x \land e \notin A_i\right) \equationexplanation{tính chất phân phối của lượng từ $\forall$ với phép $\land$}; \\
   e \in x \setminus \left(\bigcup_{i \in I}A_i\right) &\iff \forall i \in I, e \in x \setminus A_i; \\
   e \in x \setminus \left(\bigcup_{i \in I}A_i\right) &\iff e \in \bigcap_{i \in I} \left(x \setminus A_i\right).
\end{align*}

Vậy, $x \setminus \left(\bigcup_{i \in I}A_i\right) = \bigcap_{i \in I} \left(x \setminus A_i\right)$ trong mọi trường hợp. Chúng ta có điều phải chứng minh.

\stepSubexercise Với mọi $e$, \begin{equation*}
   e \in x \setminus \left(\bigcap_{i \in I}A_i\right) \iff e \in x \land e \notin \bigcap_{i \in I}A_i.
\end{equation*}
Từ $e \in x$, có $e \notin \bigcap_{i \in I}A_i$ khi và chỉ khi tồn tại $i \in I$ để $e \notin A_i$. Với giá trị $i$ đó, $e \in x \land e \notin A_i$ hay $e \in x \setminus A_i$, cho nên \begin{equation*}
   \exists i \in I, e \in x \setminus A_i,
\end{equation*}
có nghĩa là $e \in x \setminus \left(\bigcap_{i \in I}A_i\right) \implies e \in \bigcup_{i \in I} \left(x \setminus A_i\right)$. Ngược lại, nếu $e \in \bigcup_{i \in I} \left(x \setminus A_i\right)$, theo định nghĩa hợp của họ phần tử, tồn tại $i \in I$ để $e \in x \setminus A_i$. Suy ra, $e \in x \land e \notin A_i$, dẫn đến $\overline{\forall i \in I, e \in A_i}$. Vì thế, kéo theo được $e \in x \setminus \left(\bigcap_{i \in I}A_i\right)$. 

Từ lập luận hai chiều, $e \in x \setminus \left(\bigcap_{i \in I}A_i\right) \iff e \in \bigcup_{i \in I} \left(x \setminus A_i\right)$. Điều phải chứng minh.

\stepSubexercise Với mỗi $i$ trong $I$, $A_i \cap M$ là xác định duy nhất. Ngoài ra, có $\forall i \in I, A_i \cap M \subseteq x$ nên $\left(A_i \cap M\right)_{i \in I}$ là một họ phần tử của $x$.

Với mọi $e$, \begin{align*}
   e \in \left(\bigcap_{i \in I} A_i\right) \cap M &\iff e \in \bigcap_{i \in I} A_i \land e \in M; \\
   e \in \left(\bigcap_{i \in I} A_i\right) \cap M &\iff e \in x \land \forall i \in I, e \in A_i \land e \in M; \\
   e \in \left(\bigcap_{i \in I} A_i\right) \cap M &\iff e \in x \land \forall i \in I, \left(e \in A_i \land e \in M\right); \\
   e \in \left(\bigcap_{i \in I} A_i\right) \cap M &\iff e \in \bigcap_{i \in I} \left(A_i \cap M\right).
\end{align*}
Vì vậy, $\left(\bigcap_{i \in I} A_i\right) \cap M = \bigcap_{i \in I} \left(A_i \cap M\right)$. Điều phải chứng minh.

\stepSubexercise Với mỗi $i$ trong $I$, chúng ta xác định duy nhất một tập hợp $A_i \cup M$. Ngoài ra, $\forall i \in I, A_i \cup M \subseteq x$ nên $\left(A_i \cup M\right)_{i \in I}$ là một họ phần tử của $x$.

Các lập luận dưới đây đã giả thiết rằng $I \neq \varnothing$. Xét một $e \in \left(\bigcup_{i \in I}A_i\right) \cup M$, điều này tương đương với $e \in \bigcup_{i \in I}A_i \lor e \in M$. Nếu \coloringCrimson{$e \in M$}, do $I \neq \varnothing$ nên tồn tại $i \in I$ để xác định được $A_i$, kéo theo $e \in A_i \cup M$. Do đó, $e \in \bigcup_{i \in I} (A_i \cup M)$. Còn nếu \coloringDodgerblue{$e \notin M$}, khi đó phải có $e \in \bigcup_{i \in I}A_i$, có nghĩa là $\exists i \in I, e \in A_i$. Với giá trị $i$ đó, cũng sẽ có $e \in A_i \cup M$. Vì thế, $e \in \bigcup_{i \in I}(A_i \cup M)$.
Nhìn chung, có được 
\begin{equation*}
   \forall e \in \left(\bigcup_{i \in I}A_i\right) \cup M, e \in \bigcup_{i \in I}(A_i \cup M).
\end{equation*}

Ngược lại, giả sử $e \in \bigcup_{i \in I}(A_i \cup M)$. Tồn tại $i \in I$ sao cho $e \in A_i \cup M$. Khi này, $e \in A_i$ hoặc $e \in M$. Nếu $e \in A_i$ thì $e \in \bigcup_{i \in I}A_i$. Do đó, luôn có $e$ thuộc $\bigcup_{i \in I}A_i \cup M$. Vì vậy,
\begin{equation*}
   \forall e \in \bigcup_{i \in I}(A_i \cup M), e \in \bigcup_{i \in I}A_i \cup M.
\end{equation*}

Kết hợp lại, có $\forall e, \left(e \in \bigcup_{i \in I}(A_i \cup M) \iff e \in \bigcup_{i \in I}A_i \cup M\right)$ khi $I \neq \varnothing$, dẫn đến $I \neq \varnothing \implies \bigcup_{i \in I}(A_i \cup M) = \bigcup_{i \in I}A_i \cup M$. Điều phải chứng minh.

\stepSubexercise Với mọi $e$,
\begin{align*}
   e \in \left(\bigcup_{i \in I}A_i\right) \cap M &\iff e \in \left(\bigcup_{i \in I}A_i\right) \land e \in M; \\
   e \in \left(\bigcup_{i \in I}A_i\right) \cap M &\iff \exists i \in I, e \in A_i \land e \in M; \\
   e \in \left(\bigcup_{i \in I}A_i\right) \cap M &\iff j \in I \land e \in A_j \land e \in M \equationexplanation{cụ thể hóa $\exists i$ bởi $j$}; \\
   e \in \left(\bigcup_{i \in I}A_i\right) \cap M &\iff \exists i, (i \in I \land e \in A_i \land e \in M) \equationexplanation{tổng quát hóa}; \\
   e \in \left(\bigcup_{i \in I}A_i\right) \cap M &\iff \exists i \in I, e \in A_i \cap M; \\
   e \in \left(\bigcup_{i \in I}A_i\right) \cap M &\iff e \in \bigcup_{i \in I}(A_i \cap M).
\end{align*}
Vì thế, $\left(\bigcup_{i \in I}A_i\right) \cap M = \bigcup_{i \in I}(A_i \cap M)$. Điều phải chứng minh.

\stepSubexercise Với mọi $e$,
\begin{align*}
   e \in \left(\bigcap_{i \in I} A_i\right) \cup M &\iff e \in \left(\bigcap_{i \in I} A_i\right) \lor e \in M; \\
   e \in \left(\bigcap_{i \in I} A_i\right) \cup M &\iff \forall i \in I, e \in A_i \lor e \in M; \\
   e \in \left(\bigcap_{i \in I} A_i\right) \cup M &\iff \forall i, (i \in I \implies e \in A_i) \lor e \in M; \\
   e \in \left(\bigcap_{i \in I} A_i\right) \cup M &\iff \forall i , \left((i \in I \implies e \in A_i) \lor e \in M\right) \equationexplanation{cụ thể hóa và tồn tại hóa $\forall i$}; \\
   e \in \left(\bigcap_{i \in I} A_i\right) \cup M &\iff \forall i , \left(i \in I \implies e \in A_i \lor e \in M\right) \equationexplanation{$(P \implies Q) \lor R \iff (P \implies Q \lor R)$}; \\
   e \in \left(\bigcap_{i \in I} A_i\right) \cup M &\iff \forall i \in I, e \in A_i \cup M; \\
   e \in \left(\bigcap_{i \in I} A_i\right) \cup M &\iff e \in \bigcup_{i \in I}(A_i \cup M).
\end{align*}
Do vậy, $\left(\bigcap_{i \in I} A_i\right) \cup M = \bigcup_{i \in I}(A_i \cup M)$. Chúng ta có điều phải chứng minh.

\stepSubexercise Giả thiết rằng $\forall i \in I, A_i \subseteq B_i$. Với mọi $e$, nếu $e \in \bigcup_{i \in I}A_i$, tồn tại $i \in I$ để $e \in A_i$. Do $\forall i \in I, A_i \subseteq B_i$, với giá trị $i$ đó, $e \in B_i$, dẫn đến $\exists i \in I, e \in B_i$ hay $e \in \bigcup_{i \in I}B_i$. Qua đó, $\bigcup_{i \in I}A_i \subseteq \bigcup_{i \in I}B_i$.

Nếu $e \in \bigcap_{i \in I}A_i$, $\forall i \in I, e \in A_i$. Kết hợp với $\forall i \in I, A_i \subseteq B_i$, chúng ta có $\forall i \in I, e \in B_i$, tức là $e \in \bigcap_{i \in I}B_i$. Vì lí do đó, $\bigcap_{i \in I}A_i \subseteq \bigcap_{i \in I}B_i$.

Tổng hợp lại, từ $\forall i \in I, A_i \subseteq B_i$, chúng ta suy ra được $\begin{cases}
   \bigcup_{i \in I} A_i \subseteq \bigcup_{i \in I}B_i \\
   \bigcap_{i \in I} A_i \subseteq \bigcap_{i \in I}B_i \\
\end{cases}$, điều phải chứng minh.

\exercise Cho tập hợp $x$, $A$ là ánh xạ từ $I$ đến $\mathfrak{P}(x)$, và $B$ là ánh xạ từ $J$ đến $\mathfrak{P}(x)$. Chứng minh rằng
\begin{multicols}{2}
   \begin{enumerate}
      \item $\left(\bigcup_{i \in I}A_i\right) \cap \left(\bigcup_{j \in J}B_j\right) = \bigcup_{(i; j) \in I \times J}\left(A_i \cap B_j\right)$;
      \item $\left(\bigcap_{i \in I}A_i\right) \cup \left(\bigcap_{j \in J}B_j\right) = \bigcap_{(i; j) \in I \times J}\left(A_i \cup B_j\right)$.
   \end{enumerate}
\end{multicols}

